% !TEX TS-program = XeLaTeX
% use the following command:
% all document files must be coded in UTF-8
\documentclass[spanish]{textolivre}
% build HTML with: make4ht -e build.lua -c textolivre.cfg -x -u article "fn-in,svg,pic-align"

\journalname{Texto Livre}
\thevolume{19}
%\thenumber{1} % old template
\theyear{2026}
\receiveddate{\DTMdisplaydate{2026}{3}{13}{-1}} % YYYY MM DD
\accepteddate{\DTMdisplaydate{2026}{3}{19}{-1}}
\publisheddate{\DTMdisplaydate{2026}{5}{29}{-1}}
\corrauthor{Laura Ramírez}
\articledoi{10.1590/1983-3652.2026.64997}
%\articleid{NNNN} % if the article ID is not the last 5 numbers of its DOI, provide it using \articleid{} commmand 
% list of available sesscions in the journal: articles, dossier, reports, essays, reviews, interviews, editorial
\articlesessionname{reviews}
\runningauthor{Rubbo y Ramírez} 
%\editorname{Leonardo Araújo} % old template
\sectioneditorname{Daniervelin Pereira~\orcid{0000-0003-1861-3609}}
\layouteditorname{Leonardo Araújo~\orcid{0000-0003-3884-2177}}

\title{Reseña de \textit{Inteligencia artificial: 26 preguntas y respuestas}}
\othertitle{Review of \textit{Inteligencia artificial: 26 preguntas y respuestas}}
\othertitle{Resenha de \textit{Inteligencia artificial: 26 preguntas y respuestas}}
% if there is a third language title, add here:
%\othertitle{Artikelvorlage zur Einreichung beim Texto Livre Journal}

\author[1,2]{Yamila Rubbo~\orcid{0000-0002-1531-7073}\thanks{Email: \href{mailto:yamila.rubbo@conicet.gov.ar}{yamila.rubbo@conicet.gov.ar}}}
\author[1,3]{Laura Ramírez~\orcid{0000-0002-4511-9801}\thanks{Email: \href{mailto:ramirezlaura91@hotmail.com}{ramirezlaura91@hotmail.com}}}
\affil[1]{Consejo Nacional de Investigaciones Científicas y Técnicas (CONICET), Buenos Aires, Argentina.}
\affil[2]{Universidad Nacional de Tres de Febrero, Instituto de Estudios Filológicos Latinoamericanos
Ana María Barrenechea, Buenos Aires, Argentina.}
\affil[3]{Centro Interdisciplinario de investigación de psicología, matemática y experimental (CIIPME). Ciudad Autónoma de Buenos Aires, Argentina.}


\addbibresource{article.bib}
% use biber instead of bibtex
% $ biber article

% used to create dummy text for the template file
\definecolor{dark-gray}{gray}{0.35} % color used to display dummy texts
\usepackage{lipsum}
\SetLipsumParListSurrounders{\colorlet{oldcolor}{.}\color{dark-gray}}{\color{oldcolor}}

% used here only to provide the XeLaTeX and BibTeX logos
\usepackage{hologo}

% if you use multirows in a table, include the multirow package
\usepackage{multirow}

% provides sidewaysfigure environment
\usepackage{rotating}

\usepackage{endnotes} % Inclui o pacote endnotes

% CUSTOM EPIGRAPH - BEGIN 
%%% https://tex.stackexchange.com/questions/193178/specific-epigraph-style
\usepackage{epigraph}
\renewcommand\textflush{flushright}
\makeatletter
\newlength\epitextskip
\pretocmd{\@epitext}{\em}{}{}
\apptocmd{\@epitext}{\em}{}{}
\patchcmd{\epigraph}{\@epitext{#1}\\}{\@epitext{#1}\\[\epitextskip]}{}{}
\makeatother
\setlength\epigraphrule{0pt}
\setlength\epitextskip{0.5ex}
\setlength\epigraphwidth{.7\textwidth}
% CUSTOM EPIGRAPH - END

% to use IPA symbols in unicode add
%\usepackage{fontspec}
%\newfontfamily\ipafont{CMU Serif}
%\newcommand{\ipa}[1]{{\ipafont #1}}
% and in the text you may use the \ipa{...} command passing the symbols in unicode

% LANGUAGE - BEGIN
% ARABIC
% for languages that use special fonts, you must provide the typeface that will be used
% \setotherlanguage{arabic}
% \newfontfamily\arabicfont[Script=Arabic]{Amiri}
% \newfontfamily\arabicfontsf[Script=Arabic]{Amiri}
% \newfontfamily\arabicfonttt[Script=Arabic]{Amiri}
%
% in the article, to add arabic text use: \textlang{arabic}{ ... }
%
% RUSSIAN
% for russian text we also need to define fonts with support for Cyrillic script
% \usepackage{fontspec}
% \setotherlanguage{russian}
% \newfontfamily\cyrillicfont{Times New Roman}
% \newfontfamily\cyrillicfontsf{Times New Roman}[Script=Cyrillic]
% \newfontfamily\cyrillicfonttt{Times New Roman}[Script=Cyrillic]
%
% in the text use \begin{russian} ... \end{russian}
% LANGUAGE - END

% EMOJIS - BEGIN
% to use emoticons in your manuscript
% https://stackoverflow.com/questions/190145/how-to-insert-emoticons-in-latex/57076064
% using font Symbola, which has full support
% the font may be downloaded at:
% https://dn-works.com/ufas/
% add to preamble:
% \newfontfamily\Symbola{Symbola}
% in the text use:
% {\Symbola }
% EMOJIS - END

% LABEL REFERENCE TO DESCRIPTIVE LIST - BEGIN
% reference itens in a descriptive list using their labels instead of numbers
% insert the code below in the preambule:
%\makeatletter
%\let\orgdescriptionlabel\descriptionlabel
%\renewcommand*{\descriptionlabel}[1]{%
%  \let\orglabel\label
%  \let\label\@gobble
%  \phantomsection
%  \edef\@currentlabel{#1\unskip}%
%  \let\label\orglabel
%  \orgdescriptionlabel{#1}%
%}
%\makeatother
%
% in your document, use as illustraded here:
%\begin{description}
%  \item[first\label{itm1}] this is only an example;
%  % ...  add more items
%\end{description}
% LABEL REFERENCE TO DESCRIPTIVE LIST - END


% add line numbers for submission
%\usepackage{lineno}
%\linenumbers

\begin{document}
\maketitle

\begin{figure}[htbp]
\centering
\begin{minipage}{.5\textwidth}
\includegraphics[width=\textwidth]{Fig1.png}
 \caption*{\fullcite{alvarez2025inteligencia}}
 %\caption{\fullcite{nascimento_formacao_2019}.}
 \label{fig01}
\end{minipage}
\end{figure}

\begin{polyabstract}
\begin{abstract}
El presente trabajo es una reseña del libro Inteligencia artificial: 26 preguntas y respuestas \cite{alvarez2025inteligencia}, que se propone brindar conceptos, herramientas y desafíos para incorporar la inteligencia artificial (IA) en los distintos niveles educativos. La reseña presenta la organización de la obra, considerando tres apartados que analizan la práctica docente, la propuesta de actividades concretas y la apertura de interrogantes sobre las condiciones para incluir la tecnología en las aulas. Los resultados de la lectura crítica indican que este libro ofrece una mirada proactiva y reflexiva, desde la cual las autoras defienden el rol docente como mediador humano y crítico. El libro se destaca por un recorrido minucioso y extensivo de distintas herramientas digitales, pero también por la incorporación de secuencias didácticas realizadas por docentes que llevan a la reflexión sobre los desafíos de estas incorporaciones de la IA en el aula. Se concluye que el aporte del libro radica en la provisión de herramientas prácticas y en la convocatoria a las escuelas para liderar la reflexión sobre qué educación se desea construir en la era digital. 

\keywords{Inteligencia artificial \sep Educación \sep Tecnología educacional \sep Tecnología de la información}
\end{abstract}

\begin{portuguese}
\begin{abstract}
O presente trabalho é uma resenha do livro Inteligencia artificial: 26 preguntas y respuestas \cite{alvarez2025inteligencia}, que se propõe a fornecer conceitos, ferramentas e desafios para incorporar a inteligência artificial (IA) nos diferentes níveis educacionais. A resenha apresenta a organização da obra, considerando três seções que analisam a prática docente, a proposta de atividades concretas e o levantamento de questões sobre as condições para incluir a tecnologia nas salas. Os resultados da leitura crítica indicam que este livro oferece uma visão proativa e reflexiva, a partir da qual as autoras defendem o papel docente como mediador humano e crítico. Conclui-se que a contribuição do livro reside no fornecimento de ferramentas práticas e no apelo às escolas para liderarem a reflexão sobre que tipo de educação se deseja construir na era digital. 

\keywords{Inteligência artificial \sep Educação \sep Tecnologia educacional \sep Tecnologia da informação}
\end{abstract}
\end{portuguese}

\begin{english}
\begin{abstract}
This work is a review of the book Inteligencia artificial: 26 preguntas y respuestas \cite{alvarez2025inteligencia}. It aims to present concepts, tools, and challenges for incorporating artificial intelligence (AI) into the educational system. The review presents the organisation of the book considering its three sections. These analyse teaching practices, suggest specific activities, and ask questions about the conditions for using and adding AI technology in classrooms. The results of the critical reading indicate that this book provides a proactive and thoughtful point of view, from which the authors support the teachers’ role as human and critical mediators. The book stands out for its thorough examination of different digital tools, but also for the inclusion of teaching sequences designed by educators. These lead to a critical evaluation about the challenges of incorporating AI into educational environments. It is concluded that the contribution of the book makes a significant contribution by providing practical tools. Furthermore, it is a call for educational institutions to consider and take the lead on the characteristics of the education they aspire to build in the digital age.

\keywords{Artificial Intelligence \sep Education \sep Educational Technology \sep Information technology}
\end{abstract}
\end{english}
% if there is another abstract, insert it here using the same scheme
\end{polyabstract}

\section{Introdución}\label{sec-intro}
\citetitle{alvarez2025inteligencia} fue publicado en Argentina por la Editorial El Ateneo en su colección Ateneo Aula. Esta colección se orienta a repensar los procesos de aprendizaje escolar integrando el quehacer docente con la investigación científica. Las autoras, que se incorporan a esta propuesta de manera consistente, son las Dras. Guadalupe Alvarez y Lourdes Morán, investigadoras del CONICET e integrantes del Instituto del Desarrollo Humano, Universidad Nacional General Sarmiento (Buenos Aires, Argentina), y la Dra. Hilda Difabio de Anglat, también investigadora del CONICET e integrante del Centro de Investigaciones Cuyo “Dr. Abelardo Pithod” (Mendoza, Argentina). Todas forman parte del Grupo de Investigación en Pedagogías de la Lectura y la Escritura en la Cultura Digital, abocado al estudio sobre lecturas, escrituras y pedagogías en el mundo tecnológico actual. 

El libro está organizado en una introducción y tres apartados, junto con un índice de marcas mencionadas, bibliografía ampliatoria y referencias. Siguiendo la propuesta de la colección, cada apartado se subdivide en preguntas y sus respectivas respuestas. El primer apartado, “Aproximación a la inteligencia artificial en educación”, desarrolla conceptos e información clave para repensar distintas herramientas de inteligencia artificial (IA) en el ámbito educativo. El segundo, “Prácticas para integrar la IA según el nivel educativo y asignatura”, aborda posibles implementaciones para los distintos niveles y áreas desde una perspectiva conceptual y práctica. El tercer apartado, “Transformaciones de los roles de estudiante y docente”, problematiza desafíos en la formación docente, en el empleo que los/as estudiantes hacen de la IA y cómo abordarlos. El índice de marcas consta de un listado de las empresas y las herramientas digitales mencionadas en el libro. Por último, presenta bibliografía ampliatoria. Esta cuenta con sus títulos y un breve resumen de cada libro sugerido que acompaña a los/as lectores/as para profundizar el tema.

\section{Desarrollo y análisis crítico}\label{sec-normas}
El primer apartado, “Aproximación a la inteligencia artificial en educación”, enmarca el surgimiento de las IA y explica cómo funcionan para que el/la lector/a pueda comprender cabalmente sus posibilidades y limitaciones para la educación, la docencia y la enseñanza. Se introduce la historia de las IA y se presentan las específicas —aquellas creadas con objetivos específicos—. Luego explican cómo estas se desarrollan y distinguen las capacidades de aprendizaje de las IA y de las personas. Esta distinción se centra en la comprensión y la experiencia: mientras las IA pueden procesar grandes volúmenes de datos e identificar patrones basándose en reglas, carecen de una comprensión real de sus acciones. Por el contrario, las personas ponemos en juego nuestra lógica, emociones e intuición para dar sentido a la experiencia. Las siguientes preguntas exploran qué IA pueden emplearse en las aulas y brindan estrategias para hacerlo de forma adecuada. Se detienen específicamente en el uso de \textit{prompts} (instrucciones) y en qué información brindar para obtener respuestas ajustadas y precisas. En la pregunta 5 —pese al tono generalmente optimista del libro—, señalan múltiples limitaciones vinculadas con el acceso dispar a la tecnología, los sesgos muchas veces no visibilizados de las IA, el uso pasivo de los/as estudiantes y la formación docente, entre otros. Estos desafíos dialogan con los miedos docentes identificados por la literatura en torno al impacto de la IA en el pensamiento crítico, la creatividad y la integridad académica \cite{sartorharada2026professores}. En la pregunta 6, las autoras distinguen entre enseñar para, sobre y con la IA como perspectivas entrelazadas que requieren indefectiblemente de formación que acompañe su incorporación en las aulas. La pregunta 7 focaliza en cómo las herramientas digitales pueden servir para personalizar el aprendizaje a contextos diversos (socioculturales, académicos, económicos, etc.). Las últimas preguntas relevan cómo emplear la IA para producir materiales y evaluar trabajos, dos tareas que suelen agudizar la sobrecarga laboral docente. No solo se desglosan posibles herramientas, sino que también —como harán a lo largo de todo el libro— enfatizan la indispensabilidad del/la docente. En efecto, es este/a quien vuelve crítica y humana una tarea que realiza la IA. Este apartado brinda a los/as lectores/as información básica pero necesaria para que las IA, en lugar de operar como cajas negras, puedan ser verdaderas herramientas de trabajo. En síntesis, se responde a uno de los temores docentes —la falta de conocimiento sobre cómo operan las IA— al tiempo que se brindan nociones para repensar las propias prácticas.

El segundo apartado, “Prácticas para integrar la IA según el nivel educativo y asignatura”, indaga cómo puede emplearse la IA en distintos niveles educativos y áreas del conocimiento, desglosando las particularidades y potencialidades que revisten. En el nivel inicial, prioriza el juego; en el nivel primario, enfoca en Lengua y Matemáticas; en el nivel secundario, se incorporan además las áreas de Literatura, Ciencias Sociales, y Ciencias Naturales; en el nivel superior, se centra específicamente en la alfabetización académica. Las Lenguas Extranjeras se abordan en los niveles primario y secundario. Cada respuesta menciona tecnologías específicas, tipo de herramienta, función y características principales. Además, brindan ejemplos de posibles actividades en abstracto, así como de propuestas concretas desarrolladas e implementadas por docentes.

En el nivel inicial, se destaca el juego como motor de la enseñanza. Se formula un acompañamiento en la elaboración de propuestas para que la IA opere como una herramienta auxiliar. En el caso de las infancias, se destaca que el uso no puede ser autónomo ni en soledad. En el nivel primario, las propuestas se orientan tanto al uso por parte de los/as docentes —por ejemplo para personalizar sus actividades—, como de los/as estudiantes —con foco en usos críticos y regulados—. En el nivel secundario, la voz de los/as docentes ingresa de forma más clara. Se plasman usos actuales que los/as adolescentes hacen de la tecnología —y sus implicancias en el aula— y propuestas didácticas concretas en que los/as docentes enseñan a utilizar la IA. La evaluación deja de estar solamente en el resultado y pasa a estar en el proceso. En el nivel superior, el libro se estructura en dos preguntas. La primera se centra en cómo las herramientas digitales pueden asistir a los/as estudiantes en distintas tareas —búsquedas bibliográficas, elaboración de resúmenes, comparación de fuentes y su registro—. La segunda se enfoca en la elaboración de textos y la importancia de la formación para usar IA. En síntesis, las preguntas de este apartado atienden a quiénes son los/as sujetos de enseñanza y sus particularidades, así como cuáles son sus objetos de aprendizaje. Mientras en los niveles básicos las IA se presentan como acompañamiento de la tarea docente, con un acceso restringido y fuertemente mediado, en los niveles superiores se centran en los usos, sus posibilidades y limitaciones. En todas las respuestas se enfatiza fuertemente el rol mediador y crítico del docente, considerando que las nuevas generaciones no necesariamente sabrán emplearlas adecuada y críticamente sin una instrucción explícita previa.

El último apartado, “Transformaciones de los roles de estudiante y docente”, profundiza en la reconfiguración de los roles educativos. En la pregunta 21, las autoras presentan principios guía para incorporar IA en la educación para que sea una herramienta activa que promueva el desarrollo de las capacidades de los/as estudiantes. Se centran en las herramientas que pueden favorecer dos aspectos complejos: el trabajo en equipos y el diseño de propuestas focalizadas en la resolución de problemas. La pregunta 22 presenta herramientas adecuadas al desarrollo cognitivo de los distintos grupos etarios. Las preguntas 23 y 24 dan cuenta del problema del uso pasivo de las IA por parte de los/as estudiantes e indagan en estrategias para el desarrollo de competencias digitales. Para ello, ejemplifican con propuestas realizadas por docentes para atender a estos desafíos. La pregunta 25 refiere específicamente a un desafío aún pendiente: la formación docente en el uso de IA. Sugieren contemplar tres ejes básicos: explicar conceptos esenciales sobre IA, brindar habilidades pedagógicas para la IA y abordar competencias para la gestión de la IA en el aula. Por último, la pregunta 26 reflexiona sobre cómo la IA puede transformar el rol docente en el aula. Independientemente de destacar la tarea docente como una labor fuertemente humana y consecuentemente irremplazable por una máquina, las autoras enfatizan la necesidad de que este cambio no sea algo que acontece en las aulas, sino producto de una reflexión acerca de qué educación se quiere construir. El libro concluye con una invitación a no pensar a la escuela como un receptor pasivo de la IA, sino desde su rol reflexivo, a centrar el eje en cómo estas herramientas pueden ayudar a construir la educación que queremos.

\section{Conclusión}\label{sec-conduta}
\citetitle{alvarez2025inteligencia} presenta posibles usos de distintas herramientas digitales desde una mirada proactiva, que llama a los/as docentes a repensar su vínculos con las IA. Con una escritura clara, amena y didáctica, ejemplos concretos y cuadros que sintetizan información compleja de forma accesible para un público amplio, este libro resulta una puerta de entrada al empleo de la IA en la educación. Su estructura es ordenada y adecuada para docentes que no están familiarizados con la IA generativa. En este sentido, brinda herramientas prácticas con un amplio repertorio de IA disponibles y sus posibles funciones. Resulta especialmente destacable la profunda indagación que han realizado las autoras sobre la variedad de herramientas presentadas y la preocupación para pensar cómo cada una de ellas podría emplearse en los distintos campos académicos y niveles escolares. En este sentido, el libro da evidencia del relevamiento, análisis y sistematización realizadas por las autoras en investigaciones previas (ver \textcite{moran2025analisis}). 

Otro aporte destacable es el diálogo con docentes en ejercicio. Los apartados en los cuales se relevan y presentan secuencias didácticas constituyen un aporte sumamente enriquecedor, dado que se detallan la propuesta, cómo fue recibida y realizada por los/as estudiantes, y los ajustes que los/as docentes realizaron en función de las dificultades surgidas. Por ejemplo, frente al desafío de cómo evaluar trabajos realizados con IA, estas secuencias muestran cómo el prompt puede servir para dar cuenta de los saberes que los/as estudiantes ponen en juego.

En nuestro contexto actual de uso incremental de las IA en la educación, los/as docentes expresan variadas preocupaciones sobre los desafíos de la IA, a la vez que manifiestan optimismo sobre su uso en el aula \cite{almutawah2026aied,sartorharada2026professores}. Este libro da cuenta de esas inquietudes y proporciona respuestas prácticas para implementar. No obstante, destaca que el uso de la IA debe atender siempre a las características de los sujetos y a la premisa de que la labor docente es intrínsecamente humana. En este sentido, quedan pendientes cuestiones acuciantes sobre el uso de la IA, especialmente en educación: ¿qué sucede con la falta de equidad para el acceso a la tecnología? \cite{sartorharada2026professores}, ¿cómo capacitar en un uso reflexivo de las IA, es decir, una alfabetización crítica algorítmica que consista no solo en acercar conocimiento utilitario de las herramientas digitales, sino también en cómo formar docentes y estudiantes desde una perspectiva ética y política? \cite{gonsales2023ianeducacao,picao2024iaeducacao,sartorharada2026professores}. Ante un uso de la IA todavía incipiente y muchas veces acrítico, Inteligencia artificial: 26 preguntas y respuestas nos interroga sobre el impacto a largo plazo en los roles educativos y convoca a que sea la escuela la que decida qué educación construir con la IA. 


\printbibliography\label{sec-bib}
% if the text is not in Portuguese, it might be necessary to use the code below instead to print the correct ABNT abbreviations [s.n.], [s.l.]
%\begin{portuguese}
%\printbibliography[title={Bibliography}]
%\end{portuguese}

%full list: conceptualization,datacuration,formalanalysis,funding,investigation,methodology,projadm,resources,software,supervision,validation,visualization,writing,review
\begin{contributors}[sec-contributors]
\authorcontribution{Yamila Rubbo}[conceptualization,writing,review]
\authorcontribution{Laura Ramírez}[conceptualization,writing,review]
\end{contributors}

\begin{dataavailability}
\txtdataavailability{nodata} % options: dataavailable, dataonly, databody, datanotav, nodata
\end{dataavailability}

\end{document}

